\chapter{Ecuaciones diofánticas}

\section{Repaso: máximo común divisor, Euclides y Bézout}

\begin{definicion}
Sean $a,b\in\mathbb{Z}$ no ambos nulos. El \textbf{máximo común divisor} de $a$ y
$b$, que se anota $(a:b)$, es el mayor entero positivo que divide a $a$ y a $b$.
Decimos que $a$ y $b$ son \textbf{coprimos} si $(a:b)=1$.
\end{definicion}

\begin{proposicion}[{title={Lema de Euclides}}]
Si $a=qb+r$ con $q,r\in\mathbb{Z}$, entonces $(a:b)=(b:r)$.
\end{proposicion}

\begin{demostracion}
Si $x$ divide a $a$ y a $b$, entonces divide a $r=a-qb$. Recíprocamente, si $x$
divide a $b$ y a $r$, entonces divide a $a=qb+r$. Entonces $a,b$ y $b,r$ tienen
los mismos divisores comunes, y por lo tanto el mismo máximo.
\end{demostracion}

\begin{nota}
\textbf{Algoritmo de Euclides.} Se dividen sucesivamente: $a=q_1b+r_1$,
$b=q_2r_1+r_2$, $r_1=q_3r_2+r_3$, \dots, hasta que el resto sea $0$. Por el lema
de Euclides, el último resto no nulo es $(a:b)$.
\end{nota}

\begin{proposicion}[{title={Lema de Bézout}}]
Sean $a,b\in\mathbb{Z}$ no ambos nulos y $d=(a:b)$. Existen $u,v\in\mathbb{Z}$ tales
que
\[
ua+vb=d.
\]
\end{proposicion}

\begin{demostracion}
Sea $L=\{ua+vb : u,v\in\mathbb{Z}\}\cap\mathbb{N}$. Es no vacío (contiene a $|a|$
o a $|b|$), así que tiene un mínimo $e=u_0a+v_0b$. Dividimos $a=qe+r$ con
$0\le r<e$. Entonces $r=a-qe=(1-qu_0)a+(-qv_0)b$ también es combinación de $a$ y
$b$, y como $r<e$ y $e$ es el mínimo positivo, debe ser $r=0$. Luego $e\mid a$, y
análogamente $e\mid b$, así que $e\le d$. Por otro lado $d$ divide a $a$ y a $b$,
así que $d\mid u_0a+v_0b=e$, y entonces $d\le e$. Por lo tanto $e=d$.
\end{demostracion}

\begin{ejemplo}
Calcular $(84:30)$ y escribirlo como combinación de $84$ y $30$.
\[
84=2\cdot 30+24, \qquad 30=1\cdot 24+6, \qquad 24=4\cdot 6+0.
\]
El último resto no nulo es $6$, así que $(84:30)=6$. Despejando hacia atrás:
\[
6=30-24=30-(84-2\cdot 30)=3\cdot 30-84,
\]
es decir, $6=(-1)\cdot 84+3\cdot 30$. Chequeo: $-84+90=6$. $\checkmark$
\end{ejemplo}

\section{Ecuaciones diofánticas lineales}

\begin{definicion}
Una \textbf{ecuación diofántica lineal} es una ecuación de la forma
\[
aX+bY=c, \qquad a,b,c\in\mathbb{Z},
\]
de la que buscamos las soluciones \emph{enteras} $(x,y)\in\mathbb{Z}\times\mathbb{Z}$.
\end{definicion}

\begin{proposicion}[{title={Criterio de existencia}}]
Sean $a,b$ no ambos nulos y $d=(a:b)$. La ecuación $aX+bY=c$ tiene solución
entera si y sólo si $d\mid c$.
\end{proposicion}

\begin{demostracion}
($\Rightarrow$) Como $d\mid a$ y $d\mid b$, si $ax+by=c$ entonces $d\mid c$.

($\Leftarrow$) Si $c=dc'$, por Bézout existen $u,v$ con $ua+vb=d$, y entonces
$a(uc')+b(vc')=dc'=c$.
\end{demostracion}

\begin{ejemplo}
$3X+6Y=5$ no tiene soluciones enteras. En efecto, $3x+6y=3(x+2y)$, así que
tendría que ser $3\mid 5$, lo cual es absurdo. (Aquí $d=(3:6)=3\nmid 5$.)
\end{ejemplo}

\begin{nota}
\textbf{Método para resolver $aX+bY=c$.}
\begin{enumerate}
    \item Calcular $d=(a:b)$. Si $d\nmid c$, no hay solución. Si $d\mid c$,
    dividir toda la ecuación por $d$ (las soluciones no cambian) para obtener una
    ecuación con coeficientes coprimos.
    \item Buscar una \textbf{solución particular} $(x_0,y_0)$: a ojo, o con
    Euclides y Bézout, multiplicando por el término independiente.
    \item Resolver la \textbf{ecuación homogénea asociada} $aX+bY=0$.
    \item La solución general es $S=S_0+S_p$, es decir, (solución homogénea) $+$
    (solución particular).
    \item Chequear reemplazando en la ecuación.
\end{enumerate}
\end{nota}

\begin{ejemplo}
Hallar todas las soluciones enteras de $84X+30Y=18$.

\textbf{Paso 1.} $d=(84:30)=6\mid 18$, así que hay solución. Dividiendo por $6$:
$14X+5Y=3$.

\textbf{Paso 2.} Como $(14:5)=1$, buscamos $14u+5v=1$: $14\cdot(-1)+5\cdot 3=1$.
Multiplicando por $3$:
\[
14\cdot(-3)+5\cdot 9=3,
\]
así que $(x_0,y_0)=(-3,9)$ es solución particular.

\textbf{Paso 3.} Homogénea: $14x+5y=0$, es decir $14x=-5y$. Entonces $5\mid 14x$ y,
como $(5:14)=1$, $5\mid x$, o sea $x=-5k$. Reemplazando, $y=14k$. Luego
$S_0=\{(-5k,14k) : k\in\mathbb{Z}\}$.

\textbf{Paso 4.}
\[
S=\{(-3-5k,\ 9+14k) : k\in\mathbb{Z}\}.
\]

\textbf{Paso 5.} $14(-3-5k)+5(9+14k)=-42-70k+45+70k=3$. $\checkmark$
\end{ejemplo}

\section{Soluciones con restricciones}

\begin{ejemplo}
Se quiere pagar exactamente $31$ pesos con monedas de $3$ y de $5$ pesos. ¿De
cuántas maneras se puede hacer?

Sean $x$ la cantidad de monedas de $3$ e $y$ la de $5$: buscamos $(x,y)$ enteros
\textbf{no negativos} con $3x+5y=31$.

Como $(3:5)=1\mid 31$ hay solución entera. Es $3\cdot 2+5\cdot(-1)=1$, así que
multiplicando por $31$ una solución particular es $(62,-31)$. La solución general
es
\[
(x,y)=(62-5k,\ -31+3k), \qquad k\in\mathbb{Z}.
\]
Pedimos $x\ge 0$ e $y\ge 0$:
\[
62-5k\ge 0 \iff k\le \tfrac{62}{5}=12{,}4 \iff k\le 12,
\]
\[
-31+3k\ge 0 \iff k\ge \tfrac{31}{3}\approx 10{,}33 \iff k\ge 11.
\]
Luego $k\in\{11,12\}$, que da las soluciones $(x,y)=(7,2)$ y $(x,y)=(2,5)$.

\textbf{Respuesta:} hay exactamente dos maneras: $7$ monedas de $3$ y $2$ de $5$
($21+10=31$), o $2$ de $3$ y $5$ de $5$ ($6+25=31$).
\end{ejemplo}

\begin{nota}
Cuando se piden soluciones en $\mathbb{N}$ o no negativas, se resuelve primero la
ecuación en $\mathbb{Z}$ y después se imponen las desigualdades sobre $k$. Puede
haber ninguna, una cantidad finita, o (si no hay restricciones) infinitas
soluciones.
\end{nota}